\section{Semantics}
\label{sec:b-semantics}

The semantics of B developed in this work is \emph{denotational}:
each term is assigned a value in a set-theoretic domain using the ZFLean framework
developed in \Cref{ch:zflean}.
Denotational semantics provides the compositional mathematical objects required
to assign meanings to B terms, which is the central concern of the following chapters;
the emphasis is on \emph{what} a B expression computes, not on \emph{how} it is
evaluated---compared to, say, an operational semantics.
Concrete terms are interpreted after abstraction to PHOAS, as described in \Cref{sec:b-syntax}.

\subsection{Interpretation of B types}
\label{subsec:b-type-interp}

We first define the interpretation of B types as ZFC sets, denoted \(\interpZ{\cdot}\).

\begin{definition}[Type interpretation][B.BType.toZFSet][B/SemanticsPHOAS.html\#B.BType.toZFSet]
\label{def:b-type-interp}
The function $\interpZ{\cdot} : \mathsf{BType} \to \mathsf{ZFSet}$ is defined by:
\[
\begin{array}{rcl}
\interpZ{\intB} &=& \Zint,\\
\interpZ{\boolB} &=& \Zbool,\\
\interpZ{\setB\tau} &=& \powZ(\interpZ{\tau}),\\
\interpZ{\tau \prodB \sigma} &=& \interpZ{\tau} \timesZ \interpZ{\sigma}.
\end{array}
\]
\end{definition}

Agreeing on the fact that this set-theoretic interpretation is appropriate, we can show
that B types are inhabited.

\begin{lemma}[Nonemptiness][B.BType.toZFSet\_nonempty][B/SemanticsPHOAS.html\#B.BType.toZFSet_nonempty]
\label{thm:b-type-nonempty}
Interpretations of B types are nonempty, \ie
\[\forall \tau,\, \interpZ{\tau} \neq \emptyset\]
\end{lemma}
\begin{proof}
By structural induction on \(\tau\).
The sets \(\Zint\) and \(\Zbool\) are nonempty by construction;
\(\emptyset \in \powZ(S)\) for any \(S\);
the Cartesian product of two nonempty sets is nonempty.
\end{proof}

Nonemptiness ensures that open terms can be given canonical default values, and
that the semantic domain is inhabited at every type.

\subsection{Semantic domain}
\label{subsec:b-sem-domain}

We now define the semantic domain of B terms, \ie the type of interpretations of B terms.
As mentioned already, the B mathematical language is rooted in set theory, so the semantic domain
is naturally expressed over set-theoretic constructions.
Our interpretation of B terms slightly diverges from the standard description given in the
B-Book~\cite{abrial1996bbook} in that we constrain the semantic domain to carry intrinsic B types,
as motivated earlier in \Cref{rem:b-constraint-on-V}, so that we achieve that denotations are only
defined for \emph{well-typed} values with respect to our type system.
The reason for this design choice is to avoid set-theoretic pathologies related to ill-typed terms,
and provide baked-in typing constraints directly in the semantics, thus circumventing some
well-definedness side conditions---actually, all conditions related to type-correctness---in the
denotation clauses.
This is highly motivated by the idea of obtaining a \emph{correct-by-construction} design.

\begin{remark}
  This design choice only removes case analysis on type-correctness from the denotation clauses, but does not eliminate the need for semantic well-definedness side conditions related to partial operations, like in domain-application or cardinality, which are still present in the semantics and are discussed below.
\end{remark}

\begin{definition}[Semantic domain][B.Dom][B/SemanticsPHOAS.html\#B.Dom]
\label{def:b-dom}
The semantic domain is the following dependent sum type:
\[
  \mathsf{Dom} \triangleq
  \sum_{X : \mathsf{ZFSet}}\sum_{\tau : \mathsf{BType}}\,
  X \in \interpZ{\tau}
\]
\end{definition}

Elements in \(\mathsf{Dom}\) are therefore dependent triples, denoted \(\denval{X}[\tau]\),
containing a ZFC set \(X\), a B type \(\tau\) and a proof \(\pf{\tau}{X} : X \in \interpZ{\tau}\).

\begin{note}
  When the proof is unused or irrelevant in a given context, we simply write ``\(\irrelprf\)'' in
  such triples, mimicking Lean notations.
\end{note}

\subsection{Denotation}
\label{subsec:b-denote}

As motivated in \Cref{subsec:phoas-for-semantics}, we now use PHOAS representation as a trick to define the semantics of B.
The denotation function is therefore a partial function---due to B containing partial constructs requiring well-definedness side conditions---over PHOAS terms parameterized by the semantic domain.

\begin{definition}[Denotation][B.denote][B/Reasoning/Lemmas.html\#B.denote]
  \label{def:b-denote}
  The denotation function is defined by structural recursion over PHOAS terms as the partial function
\[
  \begin{array}{rcccc}
  \interpB{\cdot} & : & \mathsf{PHOAS.Term}\,\mathsf{Dom} & \rightharpoonup & \mathsf{Dom}\\
  & & t & \mapsto & \denval{X}[\tau]
  \end{array}
\]
\end{definition}
The equations defining the denotation function are detailed below, and written as \(\interpB{t} \triangleq \denval{X}[\tau]\) for each term constructor.
In each case, the construction of the proof \(\pf{\tau}{X}\) is detailed.

% In the clauses below, \(\throw\) is used when denotation cannot be defined due to a failure of a well-definedness side condition.

\paragraph{Literals and constants.}
The immediate consequence of the trick is that variables denote themselves.
For any variable \(v : \mathsf{Dom}\),
\[
  \interpB{v} \triangleq v
\]
For any integer \(n\) and Boolean \(b\),
\[
\interpB{n} \triangleq \denval{\zop{n}}[\intB],\quad
\interpB{b} \triangleq \denval{\zop{b}}[\boolB]
\]
where $\zop{n}$ and $\zop{b}$ are the respective ZFC encodings of integers and Booleans.
\begin{proof}
  \(\pf{\intB}{\zop{n}} : \zop{n} \in \Zint\) and \(\pf{\boolB}{\zop{b}} : \zop{b} \in \Zbool\) are trivially given by {\small\leancodeptr{B.ZFSet.mem\_ofInt\_Int}{ZFLean/Integers.html\#ZFSet.mem_ofInt_Int}} and {\small\leancodeptr{ZFBool.mem\_ofBool\_𝔹}{ZFLean/Booleans.html\#ZFSet.ZFBool.mem_ofBool_𝔹}} respectively.
\end{proof}

For built-in sets, the denotation is given by the interpretation of their ZFC counterparts:
\[
\interpB{\ZintB}          \triangleq \denval{\Zint}[\setB\intB],\quad
\interpB{\ZboolB}         \triangleq \denval{\Zbool}[\setB\boolB]
\]

\begin{proof}
\(\pf{\setB\intB}{\Zint} : \Zint \in \powZ(\Zint)\) and \(\pf{\setB\boolB}{\Zbool} : \Zbool \in \powZ(\Zbool)\) are trivially given by the fact that any set belongs to its own powerset.
\end{proof}

\paragraph{Pairs.}
Given \(\interpB{t} = \denval{X}[\tau]\) and \(\interpB{u} = \denval{Y}[\sigma]\), we define the denotation B pairs using Kuratowski pairs:
\[
  \interpB{t \mapletB u} \triangleq \denval{\pairZ{X}{Y}}[\tau \prodB \sigma]
\]

\begin{proof}
Using that \(\interpZ{\tau \prodB \sigma} \triangleq \interpZ{\tau} \timesZ \interpZ{\sigma}\),
\(\pf{\tau \prodB \sigma}{\pairZ{X}{Y}} : \pairZ{X}{Y} \in \interpZ{\tau} \timesZ \interpZ{\sigma}\) easily follows from composing {\small\leancodeptr{ZFSet.pair\_mem\_prod}{Mathlib/SetTheory/ZFC/Basic.html\#ZFSet.pair_mem_prod}} with the proofs \(\pf{\tau}{X}\) and \(\pf{\sigma}{Y}\).
\end{proof}

\paragraph{Arithmetic.}
Given \(\interpB{t} = \denval{X}[\intB]\) and \(\interpB{u} = \denval{Y}[\intB]\) (when the types do not match \(\intB\), the denotation is undefined), we define the denotation of arithmetic operations by lifting the corresponding ZFC operations:
\[
\interpB{t \addB u} \triangleq \denval{X \addZ Y}[\intB],\quad
\interpB{t \subB u} \triangleq \denval{X \subZ Y}[\intB],\quad
\interpB{t \mulB u} \triangleq \denval{X \mulZ Y}[\intB]
\]

\begin{proof}
  All three proofs follow from lemma {\small\leancodeptr{overloadBinOp\_mem}{Extra/ZFC_Extra.html\#overloadBinOp_mem}} composed with the proofs \(\pf{\intB}{X}\) and \(\pf{\intB}{Y}\) and the fact that all three operations are closed over \(\Zint\).
\end{proof}

\paragraph{Boolean operations.}
Given \(\interpB{t} = \denval{X}[\tau]\) and \(\interpB{u} = \denval{Y}[\sigma]\), we define the denotation of Boolean operations by lifting the corresponding ZFC operations, when the types match the expected ones; otherwise, the denotation is undefined.
\begin{align*}
  \intertext{When \(\tau = \intB\) and \(\sigma = \intB\),}
    \interpB{t \leB u}   &\triangleq \denval{X \leqZ Y}[\boolB]
  \intertext{When \(\tau = \boolB\),}
    \interpB{\notB t}    &\triangleq \denval{\negZ X}[\boolB]
  \intertext{When \(\tau = \boolB\) and \(\sigma = \boolB\),}
    \interpB{t \andB u}  &\triangleq \denval{X \andZ Y}[\boolB]
  \intertext{When \(\tau = \sigma\),}
    \interpB{t \eqB u}   &\triangleq \denval{X \eqZ Y}[\boolB]
  \intertext{When \(\sigma = \setB\tau\),}
    \interpB{t \inB u}   &\triangleq \denval{X \in Y}[\boolB]
\end{align*}

\begin{proof}
  All proofs follow directly from the corresponding ZFC lemmas about the respective operations, composed with the proofs \(\pf{\tau}{X}\) and \(\pf{\sigma}{Y}\).
\end{proof}

\paragraph{Set-theoretic constructors.}
Given \(\interpB{S} = \denval{S'}[\setB\tau]\) and \(\interpB{T} = \denval{T'}[\setB\sigma]\), we define the denotation of set-theoretic constructors by lifting the corresponding ZFC operations, when the types match the expected ones; otherwise, the denotation is undefined.
Those equations witness the fact that the B language is rooted in set theory.

\begin{align*}
  \interpB{S \prodB T} &\triangleq \denval{S' \timesZ T'}[\setB(\tau \prodB \sigma)]\\
  \interpB{\powB(S)}   &\triangleq \denval{\powZ(S')}[\setB(\setB\tau)]
\intertext{
  \begin{proof}
    Using that \(\interpZ{\setB(\tau \prodB \sigma)} \triangleq \powZ(\interpZ{\tau} \timesZ \interpZ{\sigma})\), we show
    \begin{align*}
      \undernote{S' \timesZ T' \in \interpZ{\setB(\tau \prodB \sigma)}}{\pf{\setB(\tau \prodB \sigma)}{S' \timesZ T'}} &\iff S' \timesZ T' \in \powZ(\interpZ{\tau} \timesZ \interpZ{\sigma})\\[-18pt]
      &\iff S' \timesZ T' \subseteq \interpZ{\tau} \timesZ \interpZ{\sigma}\\
      &\iff S' \subseteq \interpZ{\tau} \land T' \subseteq \interpZ{\sigma}\\
      &\iff \undernote{S' \in \powZ(\interpZ{\tau})}{\pf{\tau}{S'}} \land \undernote{T' \in \powZ(\interpZ{\sigma})}{\pf{\sigma}{T'}}
    \end{align*}
    Likewise, we construct the proof \(\pf{\setB(\setB\tau)}{\powZ(S')}\).
  \end{proof}
  When \(\tau = \sigma\),}
  \interpB{S \cupB T} &\triangleq \denval{S' \cupZ T'}[\setB\tau]\\
  \interpB{S \capB T} &\triangleq \denval{S' \capZ T'}[\setB\tau]
\end{align*}
\begin{proof}
  \(\pf{\setB\tau}{S' \cupZ T'} : S' \cupZ T' \in \powZ(\interpZ{\tau})\) and \(\pf{\setB\tau}{S' \capZ T'} : S' \capZ T' \in \powZ(\interpZ{\tau})\) are easily constructed using monotonicity of union and intersection: {\small\leancodeptr{union\_mono}{Extra/ZFC_Extra.html\#union_mono}}, {\small\leancodeptr{inter\_mono}{Extra/ZFC_Extra.html\#inter_mono}}.
\end{proof}

\paragraph{Partial functions.}

Partial function spaces are interpreted as sets of functional relations.
Given \(\interpB{A} = \denval{A'}[\setB\tau]\) and \(\interpB{B} = \denval{B'}[\setB\sigma]\), let
\[
  \mathcal{F} \coloneq \bigl\{ f \in \powZ(A' \timesZ B') \bigm| \leancodeptr{IsPFunc}{ZFLean/Functions.html\#ZFSet.IsPFunc}\ f\ A'\ B'  \bigr\}\]
Then, the denotation is given by
\[
  \interpB{A \pfunB B}
  =
  \denval{\mathcal{F}}[\setB\bigl(\setB(\tau \prodB \sigma)\bigr)]
\]
\begin{proof}
  Using that \textbf{partial} functions are covariant in both their domain and codomain---see lemma {\small\leancodeptr{prod\_sep\_is\_pfunc\_mem}{ZFLean/Functions.html\#ZFSet.prod_sep_is_pfunc_mem}}, we show that
\begin{align*}
  F &\subseteq \bigl\{ f \in \powZ(\interpZ{\tau} \timesZ \interpZ{\sigma}) \bigm| \leancodeptr{IsPFunc}{ZFLean/Functions.html\#ZFSet.IsPFunc}\ f\ \interpZ{\tau}\ \interpZ{\sigma}  \bigr\} & \text{(using \(\pf{\setB\tau}{A'}\) and \(\pf{\setB\sigma}{B'}\))}\\
  & \subseteq \powZ(\interpZ{\tau} \timesZ \interpZ{\sigma}) & \text{(using \leancodeptr{sep\_subset}{Mathlib/SetTheory/ZFC/Basic.html\#ZFSet.sep_subset})}
\end{align*}%
\ie that \(F \in \powZ(\powZ(\interpZ{\tau} \timesZ \interpZ{\sigma})) = \interpZ{\setB(\setB(\tau \prodB \sigma))}\).
\end{proof}

\paragraph{Partial operations.}
The semantics of partial operations is given by case analysis on the well-definedness side conditions, and is defined only when these conditions are satisfied.

Function application evaluates by graph lookup, well-defined when the function denotes a partial function and the argument is in its domain.
Given \(\interpB{f} = \denval{F}[\setB(\tau \prodB \sigma)]\) and \(\interpB{x} = \denval{X}[\tau]\), if \(\mathsf{IsPFunc}\ F\ \interpZ{\tau}\ \interpZ{\sigma}\) and \(X \in \leancodeptr{Dom}{ZFLean/Functions.html\#ZFSet.Dom}(F)\), let \(Y\) be the unique value such that \(\pairZ{X}{Y} \in F\); the denotation is then given by
\[
  \interpB{f(x)} = \denval{Y}[\sigma]
\]
\begin{proof}
  Both uniqueness and the fact that \(Y \in \interpZ{\sigma}\) follow from the definition of \(\mathsf{IsPFunc}\).
\end{proof}
 
Cardinality is defined under finiteness conditions, like minimum and maximum, which also require nonemptiness; they lift the corresponding ZFC operations.
Given \(\interpB{S} = \denval{S'}[\setB\tau]\) such that \(\leancodeptr{IsFinite}{ZFLean/Functions.html\#ZFSet.IsFinite}\ S'\), we define:
\begin{align*}
  \interpB{\cardB{S}} &\triangleq \denval{\cardZ{S'}}[\intB]
\intertext{
  \begin{proof}
    By definition, \(\cardZ{S'}\) is defined, when \(S'\) is finite, as the unique natural number \(n\), lifted as integer, such that there exists a bijection between \(S'\) and the set \(\{0, 1, \ldots, n-1\}\); in particular, \(\cardZ{S'} \in \Zint = \interpZ{\intB}\).
  \end{proof}  
  Moreover, when \(\tau = \intB\) and \(S' \neq \emptyset\),}
  \interpB{\minB(S)} &\triangleq \denval{\minZ(S')}[\intB]\\
  \interpB{\maxB(S)} &\triangleq \denval{\maxZ(S')}[\intB]
\end{align*}
\begin{proof}
  We sketch the proof for the minimum case, the maximum case being similar.
  Using the fact that the minimum of a non-empty finite set belongs to the set (see {\small\leancodeptr{Min\_mem\_of\_non\_empty\_finite}{ZFLean/Functions.html\#ZFSet.Min_mem_of_non_empty_finite}}) and \(\pf{\setB\intB}{S'} : S' \in \powZ(\Zint)\), we show that
  \[
    \minZ(S') \in S' \subseteq \Zint = \interpZ{\intB}
    \vspace{-\baselineskip}
  \]
\end{proof}

\paragraph{Binders.}

Binder constructs rely on the PHOAS representation. Set comprehension is easily interpreted as
ZFC separation, lambda abstraction as graph construction---so, an underlying separation---and
universal quantification as bounded infimum---\ie a set-level conjunction folded over the
quantification domain---in the two-element Boolean algebra.

We start with the most straightforward case of set comprehension \(\{D \mid P\}\).
Definedness assumptions are the following:
\begin{itemize}
  \item \(\interpB{D} = \denval{\mathcal{D}}[\setB\tau^{\times}]\), where \(\tau^{\times}\) is a product type \(\tau_1 \prodB \cdots \prodB \tau_n\) of the binder's arity;
  \item For any family of variables \(\vec x = (x_1, \cdots, x_n)\) of the binder's arity such that each of the \(x_i\) carries the B type \(\tau_i\), and such that \(\pairZ{x_1}{\cdots, x_n} \in \mathcal{D}\), \(P(\vec x)\) has a denotation.
\end{itemize}
We can therefore define the following scoped denotation \(\mathbb{P} : \mathsf{ZFSet} \to \mathsf{Dom}\) of the body \(P\), for any set argument \(z\), by the following case distinction:
\begin{description}
  \item[if \(z \in \mathcal{D}\):] \(z\) can be written as an \(n\)-ary set-level tuple \(\pairZ{z_1}{\cdots, z_n}\).
  It easily follows from \(\pf{\setB\tau^\times}{\mathcal{D}}\) that each \(z_i\) belongs to \(\interpZ{\tau_i}\), thereby constructing \(\pf{\tau_i}{z_i}\).
  By assumption, \(P\) therefore has a denotation for the family of variables \((\denval{z_i}[\tau_i])_{1 \leq i \leq n}\), whose first component is a set \(p_z\), \ie
  \[
    \interpB{P\bigl((\denval{z_i}[\tau_i])_{1 \leq i \leq n}\bigr)} = \denval{p_z}[\tau_{p_z}]
  \]
  
  In this case, we define the denotation of the body at \(z\) as the following equality:
  \[
    \mathbb{P}(z) : \mathsf{Prop} \coloneq p_z = \topZ
  \]
  \begin{note}
  Note that \(\tau_{p_z}\) is not used in the definition of \(\mathbb{P}\), otherwise we would require that it be equal to \(\boolB\).
  It is unnecessary at this point and the fact can be easily retrieved when we add typing assumptions.
  \end{note}
  \item[otherwise:] we define \(\mathbb{P}(z) : \mathsf{Prop} \coloneq \mathsf{False}\).
\end{description}

The denotation of the whole comprehension is then simply given by the ZFC separation \(\leancodeptr{ZFSet.sep}{Mathlib/SetTheory/ZFC/Basic.html\#ZFSet.sep}\ \mathcal{D}\ \mathbb{P}\), or with usual set-builder notation:
\[
  \interpB{\{D \mid P\}}
  \triangleq
  \denval{\{ z \in \mathcal{D} \mid \mathbb{P}(z) \}}[\setB\tau^\times][\pf{\setB\tau^{\times}}{\{ \mathcal{D} \mid \mathbb{P} \}}]
\]

\begin{proof}
  The proof is straightforward:
\[
  \{ z \in \mathcal{D} \mid \mathbb{P}(z) \} \subseteq \undernote{\mathcal{D} \subseteq \interpZ{\tau^\times}}{\pf{\setB\tau^\times}{\mathcal{D}}}
  \vspace{-\baselineskip}
\]
\end{proof}

For lambda abstraction \(\lambdaB' D \cdot E\), the result is the graph of the function described by
the body. The definition is quite similar to set comprehension, with the technicality that the return type of the body is not known at this point without typing information, so we need to compute it on the fly as we define the denotation of the body. The assumptions are also similar:
\begin{itemize}
  \item \(\interpB{D} = \denval{\mathcal{D}}[\setB\tau^{\times}]\) as before;
  \item For any family of variables \(\vec x = (x_1, \cdots, x_n)\) of the binder's arity such that each of the \(x_i\) carries the B type \(\tau_i\), and such that \(\pairZ{x_1}{\cdots, x_n} \in \mathcal{D}\), \(E(\vec x)\) has a denotation.
\end{itemize}
We define the following scoped denotation \(\mathcal{E} : \mathsf{ZFSet} \to \mathsf{Dom}\) of the body \(E\), for any set argument \(z\), by distinguishing whether \(\mathcal{D}\) is empty or not:
\begin{description}
  \item[if \(\mathcal{D} = \emptyset\):] in this case, we cannot directly compute the return type of the body, as we have no element in the domain to plug in it; however, we can still define a dummy family of variables from the domain's type, and use it to compute the return type of the body. We have no guarantee that the body has a denotation for this family of variables, and even if it does, the denotation might be ill-typed.
  \begin{note}
    Even though this last remark may seem degenerate or even pathological, it is important to note that it is not. If such an ill-formed term were to occur as sub-calculation in a denotation clause, the denotation of the whole term would simply fail since it would not satisfy well-definedness assumptions, thus remaining equi-consistent with the semantics of B.
  \end{note}
  Such a dummy family can be defined by picking a arbitrary elements \(d_i\) in each \(\interpZ{\tau_i}\), which are nonempty by \Cref{thm:b-type-nonempty}, and defining the family as \((\denval{d_i}[\tau_i])_{1 \leq i \leq n}\).
  If \(E\bigl((\denval{d_i}[\tau_i])_{1 \leq i \leq n}\bigr)\) has a denotation whose second component is a B type \(\gamma\), we define whole denotation of the abstraction as:
  \[
    \interpB{\lambdaB' D \cdot E} \triangleq \denval{\emptyset}[\setB(\setB(\tau^\times \prodB \gamma))]
  \]
  \begin{proof}
    Trivial by {\small\leancodeptr{empty\_subset}{Mathlib/SetTheory/ZFC/Basic.html\#ZFSet.empty_subset}}.
  \end{proof}
    
  \item[otherwise:] in this case, we pick an arbitrary element \(d \in \mathcal{D}\), which can be written as an \(n\)-ary set-level tuple \(\pairZ{d_1}{\cdots, d_n}\) as before, and we define the family of variables as \((\denval{d_i}[\tau_i])_{1 \leq i \leq n}\).
  By assumption, \(E\) has a denotation for this family of variables, whose second component is a B type \(\gamma\), which is the return type of the body.
  
  The scoped denotation \(\mathcal{E}\) of the body is then defined for arbitrary set arguments \(z\) similarly to the set comprehension case.
  The only difference is that \(z\) is expected to be a set-level pair of the form \(z = \pairZ{x}{y}\).
  Similarly to set comprehension, we distinguish whether \(x\) belongs to \(\mathcal{D}\):
  \begin{description}
    \item[if \(x \in \mathcal{D}\):] we can write \(x\) as an \(n\)-ary set-level tuple \(\pairZ{x_1}{\cdots, x_n}\), where each \(x_i\) belongs to \(\interpZ{\tau_i}\) as before.
    By assumption, \(E\) has a denotation for the family of variables \((\denval{x_i}[\tau_i])_{1 \leq i \leq n}\), whose first component is a set \(e_x\). We define the denotation of the body at \(z\) as the following equality:
    \[
      \mathcal{E}(z) : \mathsf{Prop} \coloneq e_x = y
    \]
    \item[otherwise:] we define \(\mathcal{E}(z) : \mathsf{Prop} \coloneq \mathsf{False}\).
  \end{description}
  The denotation of the whole abstraction is then given using \(\mathcal{E}\) by the set-level separation
  \(\Lambda \coloneq \{ z \in \powZ(\interpZ{\tau^\times} \timesZ \interpZ{\gamma}) \mid \mathcal{E}(z) \}\):
  \[
    \interpB{\lambdaB' D \cdot E} \triangleq \denval{\Lambda}[\setB(\setB(\tau^\times \prodB \gamma))]
  \]
\end{description}

\begin{proof}
  Similar to the set comprehension case.
\end{proof}

Universal quantification is interpreted as the bounded infimum of the
values obtained over the domain, taken in the two-element ZFC Boolean
algebra $\Zbool$.
The assumptions are again similar to the previous cases:
\begin{itemize}
  \item \(\interpB{D} = \denval{\mathcal{D}}[\setB\tau^{\times}]\), with \(\tau^\times = \tau_1 \times \cdots \times \tau_n\) as before;
  \item For any family of variables \(\vec x = (x_i)_{1 \leq i \leq n}\) of the binder's arity such that each of the \(x_i\) carries the B type \(\tau_i\), and such that \(\pairZ{x_1}{\cdots, x_n} \in \mathcal{D}\), \(P(\vec x)\) has a denotation.
\end{itemize}
The scoped denotation \(\mathbb{P}\) of the body \(P\) is defined exactly as in the set comprehension case.
The intuition is the following: when the domain \(D\) is finite, so is its denotation \(\mathcal{D} = \{e_1, \cdots, e_k\}\). Universal quantification over that domain is then naturally expressed as the following enumerated conjunction:
\[
  \interpB{\forallB' D \cdot\,P} = \mathbb{P}(e_1) \andZ \cdots \andZ \mathbb{P}(e_k) = \bigandZ_{i = 1}^k \mathbb{P}(e_i) = \bigandZ_{e \in \mathcal{D}} \mathbb{P}(e)
\]

The implementation differs from this equation in two negligible technicalities.
When the domain is empty, \ie \(\mathcal{D} = \emptyset\), the infimum should be the top element in the algebra, namely \(\topZ\), also expressing vacuous truth; however in Lean, this infimum is defined as quantified intersection, which is defined, in standard set theory, over non-empty families of sets. In our model, quantified intersection over the empty set is defined as the empty set ({\small\leancodeptr{ZFSet.sInter\_empty}{Mathlib/SetTheory/ZFC/Basic.html\#ZFSet.sInter_empty}}).
In order to circumvent this first technicality, we distinguish the empty domain case and return
\(\denval{\topZ}[\boolB][{\small\leancodeptr{zftrue\_mem\_𝔹}{ZFLean/Booleans.html\#ZFSet.ZFBool.zftrue_mem_𝔹}}]\)
directly in that case.
The second technicality is that our model of ZFC does not offer generalized comprehension, so we write the following quantified intersection over a separation instead: 
\[
\bigandZ_{\mathcal{D}} \mathbb{P} = \bigcap \{z \in \ZboolB \mid \exists x \in \mathcal{D},\, y = \mathbb{P}(x) \}
\]
Overall, the denotation of the whole quantification is then given by
\[
  \interpB{\forallB' D \cdot\,P}
  \triangleq
  \denval{\bigcap \{z \in \ZboolB \mid \exists x \in \mathcal{D},\, y = \mathbb{P}(x) \}}[\boolB][\pf{}{}]
\]

\begin{proof}
  The proof \(\pf{}{} : \bigcap \{z \in \ZboolB \mid \exists x \in \mathcal{D},\, y = \mathbb{P}(x) \} \in \interpZ{\boolB}\) naturally follows from the algebra being structurally closed under conjunction.
\end{proof}

\subsection{Basic properties}
\label{subsec:b-sem-properties}

\inlremark[]{This section is still in progress, in particular proofs need further polishing.}

The denotation enjoys two fundamental properties: \emph{type adequacy} and \emph{soundness}.
Type adequacy ensures that whenever a term denotes, the intrinsic B-type tag carried by its
denotation coincides with the static type assigned by the typing derivation.
Soundness then ensures that, under suitable hypotheses on the environment and on the
well-definedness of partial operations, every well-typed term does indeed denote a value of
the expected type.

Both statements presuppose that the variables appearing in a term are bound to domain
elements whose intrinsic B types match what the abstract typing context records.
Since PHOAS variables are themselves domain elements, this coherence is conveniently
captured at the PHOAS level by the following condition, which is automatically satisfied by
contexts obtained through the abstraction \(\aterm{\Gamma}{\Delta}\) of
\Cref{def:b-context-abstraction}.

\begin{definition}[Well-typed PHOAS context][B.PHOAS.WellTypedCtx][B/Reasoning/DenotationTotality.html\#B.PHOAS.WellTypedCtx]
\label{def:b-welltyped-ctx}
A PHOAS typing context \(\Gamma : \mathsf{Dom} \rightharpoonup \mathsf{BType}\) is
\emph{well-typed}, written \(\WTctx(\Gamma)\), whenever every binding agrees with the
intrinsic type tag of its key:
\[
  \forall d : \mathsf{Dom},\, \forall \tau : \mathsf{BType},\quad
  \Gamma(d) = \tau \implies d.\tau = \tau,
\]
where \(d.\tau\) is the second projection of \(d \in \mathsf{Dom}\).
\end{definition}

\begin{theorem}[Type adequacy][denote\_welltyped\_eq][B/Reasoning/Lemmas.html\#denote_welltyped_eq]
\label{thm:b-denote-type-det}
Let \(t : \mathsf{PHOAS.Term}\,\mathsf{Dom}\) be a PHOAS term well-typed of type
\(\tau_0\) under some well-typed context, and assume
\(\interpB{t} = \denval{V}[\tau]\). Then \(\tau = \tau_0\).
\end{theorem}

\begin{proof}
By structural induction on \(t\). The key observation is that, in every clause of the
denotation function, the type tag of the result is constructed from the type tags of the
subterm denotations by the very same combinator that the typing rule uses to combine the
subterm types. The induction therefore reduces to checking the cases one by one.

For the variable case \(t = v\), the denotation returns \(v\) itself, so its type tag is
\(v.\tau\). Inversion of the variable typing rule yields \(\Gamma(v) = \tau_0\), and
\(\WTctx(\Gamma)\) gives \(v.\tau = \tau_0\), hence \(\tau = \tau_0\).
For each literal or constant (\(\mathsf{int}\,n\), \(\mathsf{bool}\,b\), \(\ZintB\),
\(\ZboolB\)), the typing rule fixes the type uniquely and the denotation tags the result
with the same type; the equality is then immediate.

For each compound non-binder constructor (pairs, arithmetic, Boolean and set-theoretic
operations, partial function space), inversion of the typing rule provides typings of the
subterms, the induction hypothesis identifies the type tags of their denotations with the
typing-derivation types, and the denotation clause assembles the result type by the same
formula as the typing rule. For example, for a pair \(t \mapletB u\), inversion yields
\(\btype{\cdot}{t}{\alpha}\) and \(\btype{\cdot}{u}{\beta}\), the induction hypothesis
matches the subterm denotation tags to \(\alpha\) and \(\beta\), and the clause sets the
result type to \(\alpha \prodB \beta\), which is exactly the typing-rule output.

For the partial operations \(\appB f\,x\), \(\cardB{S}\), \(\minB(S)\), and \(\maxB(S)\),
the denotation enters its success branch only after a side condition is verified, and that
branch tags the result with a fixed type (\(\beta\) for \(\appB\), \(\intB\) for the other
three) matching the typing-rule output.

For binders \(\{D \mid P\}\), \(\lambdaB' D \cdot E\), and \(\forallB' D \cdot P\), the
induction hypothesis on the domain identifies its type tag with the folded product
\(\setB\tau^\times\); the binder clause then constructs the result with type tag
\(\setB\tau^\times\), \(\setB(\tau^\times \prodB \gamma)\), or \(\boolB\) respectively,
matching the typing rule.
\end{proof}

In other words, a term cannot denote a value of two distinct types: the static type fully
determines the type tag of the denotation whenever the latter exists.

We now turn to existence. To formulate it, we briefly recall the well-definedness side
conditions that B-method proof obligations require for partial operations. These are the
B-side analogue of the well-definedness obligations of Atelier B and are detailed in
\Cref{sec:b-wd}; here we only need the following packaging.

\begin{definition}[Well-definedness predicate][B.PHOAS.WellDefined][B/Reasoning/WellDefined.html\#B.PHOAS.WellDefined]
\label{def:b-welldefined}
The predicate \(\WD : \mathsf{PHOAS.Term}\,\mathsf{Dom} \to \mathsf{Prop}\) is defined by
structural recursion. It holds vacuously on variables and literals, propagates conjunctively
to immediate subterms, and additionally carries the following semantic side conditions at
partial operations:
\begin{description}
  \item[\(\appB f\,x\):] writing \(\interpB{f} = \denval{F}[\setB(\tau \prodB \sigma)]\)
  and \(\interpB{x} = \denval{X}[\tau]\), the ZFC component \(F\) is a partial function
  from \(\interpZ{\tau}\) to \(\interpZ{\sigma}\), and \(X \in \domB(F)\);
  \item[\(\cardB{S}\):] writing \(\interpB{S} = \denval{S'}[\setB\alpha]\), the set
  \(S'\) is finite;
  \item[\(\minB(S),\maxB(S)\):] \(S'\) is finite and nonempty.
\end{description}
\end{definition}

The soundness theorem then asserts that well-typed, well-defined terms always denote at the
expected type. We state the result both at the PHOAS level (the form actually proved) and
at the concrete level (its intended use, obtained by composition with the typing
abstraction result of \Cref{thm:b-typing-of-abstract}).

\begin{theorem}[Denotation soundness][B.denote\_exists\_of\_typing][B/Reasoning/DenotationTotality.html\#B.denote_exists_of_typing]
\label{thm:b-denote-sound}
Let \(\Gamma : \mathsf{Dom} \rightharpoonup \mathsf{BType}\) be a PHOAS context with
\(\WTctx(\Gamma)\), and let \(t : \mathsf{PHOAS.Term}\,\mathsf{Dom}\) satisfy
\(\batype{\Gamma}{t}{\tau}\) and \(\WD(t)\). Then there exists \(V \in \interpZ{\tau}\) such
that
\[
  \interpB{t} = \denval{V}[\tau].
\]
In particular, for any concrete derivation \(\btype{\Gamma}{t}{\tau}\) and any renaming
\(\Delta\) well-formed with respect to \(\Gamma\) and covering the free variables of \(t\),
if \(\WD(\aterm{t}{\Delta})\) holds, then there exists \(V \in \interpZ{\tau}\) such that
\(\interpB{\aterm{t}{\Delta}} = \denval{V}[\tau]\).
\end{theorem}

\begin{proof}
The second statement follows from the first by composition with
\Cref{thm:b-typing-of-abstract} (\(\baop\vdash\) is preserved by abstraction) and the fact
that \(\aterm{\Gamma}{\Delta}\) is automatically well-typed in the sense of
\Cref{def:b-welltyped-ctx}, by construction.

For the PHOAS-level statement, we proceed by structural induction on the typing
derivation \(\batype{\Gamma}{t}{\tau}\). The argument splits into four groups of cases.

\paragraph{Variables and literals.}
For \(t = v\), take \(V \coloneq v.X\); the denotation clause returns \(v\) itself, and
\(V \in \interpZ{v.\tau} = \interpZ{\tau}\) follows from
\(\WTctx(\Gamma)\). For literals, the witnesses are immediate:
\(V \coloneq \zop{n} \in \Zint\) for \(\mathsf{int}\,n\),
\(V \coloneq \zop{b} \in \Zbool\) for \(\mathsf{bool}\,b\),
and \(V \coloneq \Zint \in \powZ(\Zint)\),
\(V \coloneq \Zbool \in \powZ(\Zbool)\) for \(\ZintB\) and \(\ZboolB\). In each case, the
membership proof was constructed in \Cref{subsec:b-denote}.

\paragraph{Compound non-partial constructors.}
For each of \(\mapletB\), \(\addB\), \(\subB\), \(\mulB\), \(\andB\), \(\notB\), \(\eqB\),
\(\leB\), \(\inB\), \(\powB\), \(\prodB\), \(\cupB\), \(\capB\), and \(\pfunB\), the
induction hypothesis applied to each subterm produces a witness denotation whose type tag
matches the typing-rule premise. By \Cref{thm:b-denote-type-det}, every case analysis on
subterm type tags inside the corresponding denotation clause selects the success branch.
The witness \(V\) is then constructed by applying the relevant ZFC closure lemma to the
subterm witnesses: \(\pairZ{X}{Y} \in \interpZ{\alpha} \timesZ \interpZ{\beta}\) for
maplets, closure of \(\Zint\) and \(\Zbool\) under the arithmetic and Boolean operations
for the connectives, monotonicity of \(\cupZ\), \(\capZ\), \(\timesZ\), and \(\powZ\) under
inclusion for the set constructors, and the closure of partial-function sets under the
covariant action of their domain and codomain for \(\pfunB\) (\Cref{subsec:b-denote}).

\paragraph{Partial operations.}
For \(\appB f\,x\), the induction hypothesis gives denotations
\(\interpB{f} = \denval{F}[\setB(\tau \prodB \sigma)]\) and
\(\interpB{x} = \denval{X}[\tau]\). The well-definedness hypothesis \(\WD(\appB f\,x)\)
supplies a proof that \(F\) is a partial function from \(\interpZ{\tau}\) to
\(\interpZ{\sigma}\) and a witness \(Y_0\) with \(\pairZ{X}{Y_0} \in F\); these unlock
the conditional in the denotation clause, and the partial-function property yields a
unique \(Y \in \interpZ{\sigma}\) with \(\pairZ{X}{Y} \in F\), which we take as \(V\).
For \(\cardB{S}\), \(\minB(S)\), and \(\maxB(S)\), writing
\(\interpB{S} = \denval{S'}[\setB\alpha]\), \(\WD\) supplies the finiteness (and
nonemptiness, for \(\minB\) and \(\maxB\)) hypotheses that activate the corresponding
denotation clauses; the witnesses \(\cardZ{S'}\), \(\minZ(S')\), and \(\maxZ(S')\) belong
to \(\Zint\) by their defining properties.

\paragraph{Binders.}
For \(\{D \mid P\}\), the induction hypothesis on \(D\) (which is itself a foldl of
Cartesian products of the component domains \(D_i\)) gives a denotation
\(\interpB{D} = \denval{\mathcal{D}}[\setB\tau^\times]\). By \Cref{thm:b-type-nonempty},
each \(\interpZ{\tau_i}\) is nonempty, so we can pick a canonical default tuple
\(\vec d^{\,\mathsf{def}} \in \interpZ{\tau^\times}\) with type tags exactly matching
\(\vec\alpha\); the induction hypothesis applied to \(P(\vec d^{\,\mathsf{def}})\) under the
extended context \(\Gamma[\vec v : \vec\alpha]\)---which remains well-typed since each
\(d_i^{\mathsf{def}}\) carries the intrinsic type \(\alpha_i\)---yields a Boolean
denotation, which activates the success branch of the comprehension clause. The witness is
then \(V \coloneq \{ z \in \mathcal{D} \mid \mathbb{P}(z) \}\), which lies in
\(\powZ(\interpZ{\tau^\times})\) by the separation principle.

The \(\forallB' D \cdot P\) case proceeds analogously, with two subcases: when
\(\mathcal{D} = \emptyset\), the denotation evaluates to \(\topZ \in \Zbool\); when
\(\mathcal{D} \neq \emptyset\), pick any \(d \in \mathcal{D}\) and apply the induction
hypothesis to \(P(\vec d)\) (with type tags matching \(\vec\alpha\), by membership of
\(\vec d\) in \(\mathcal{D} \subseteq \interpZ{\tau^\times}\)) to ensure the body denotes
at every tuple of the right shape; the bounded infimum then lies in \(\Zbool\) by closure
of the two-element algebra under conjunction.

The \(\lambdaB' D \cdot E\) case mirrors \(\forallB'\), with the additional twist that the
return type \(\gamma\) of the body must be retrieved before the result type
\(\setB(\tau^\times \prodB \gamma)\) can be formed. When \(\mathcal{D} = \emptyset\), the
canonical default tuple from \Cref{thm:b-type-nonempty} provides this information, and the
denotation is \(\emptyset \in \powZ(\interpZ{\tau^\times \prodB \gamma})\). When
\(\mathcal{D} \neq \emptyset\), an arbitrary witness in \(\mathcal{D}\) supplies a
representative for which the induction hypothesis returns \(\gamma\) and a denotation; the
denotation is then a subset of \(\interpZ{\tau^\times} \timesZ \interpZ{\gamma}\) by
separation.
\end{proof}

The following corollary, used freely in the correctness proofs of subsequent chapters,
records the most common form of the soundness theorem: the conclusion is restated as the
defining property of the \(\mathsf{some}\) constructor.

\begin{corollary}[Totality on well-typed, well-defined terms][B.PHOAS.denote\_isSome\_of\_typing][B/Reasoning/DenotationTotality.html\#B.PHOAS.denote_isSome_of_typing]
\label{cor:b-denote-isSome}
Under the hypotheses of \Cref{thm:b-denote-sound}, \(\interpB{t}\) is defined, \ie
\(\interpB{t} \neq \bot\).
\end{corollary}

The syntax, typing, and semantics described in this chapter form the source-language
contract used by the SMT translation and by the correctness arguments in the following
chapters.
